% !PRINTED 51
\weberpage{51}

%<<CONT
und abermals angewandt:
\weq{9}{\begin{gathered}
\mathrm{sn}\,(u+4K) = \mathrm{sn}\,u,\\
\mathrm{cn}\,(n+4K) = \mathrm{cn}\,u, % [sic] \textquotedbl{}n\textquotedbl{} printed for \textquotedbl{}u\textquotedbl{}
\\
\mathrm{dn}\,(u+4K) = \mathrm{dn}\,u.
\end{gathered}}
Es haben also die elliptischen Funktionen die Eigenschaft der \gesperrt{Periodizität} mit den trigonometrischen Funktionen gemein. Sie haben die Periode $4K$ (dn $u$ auch die Periode $2K$).

In §\,3 haben wir als Beispiel die Transformation, die Formel (1), abgeleitet:
\weq{10}{\frac{dz}{\sqrt{z(1-z)(1-\varkappa^2 z)}} = \frac{dx}{i\sqrt{x(1-x)(1-\varkappa'^2 x)}},}
wenn
\weq{11}{z = \frac{-x}{1-x}}
war. Setzt man also
\[
\frac{dx}{\sqrt{x(1-x)(1-\varkappa'^2x)}} = i\,du,
\]
so ist, wenn $x$ und $u$ zugleich verschwinden, $x = \mathrm{sn}^2(iu,\varkappa')$, $1-x = \mathrm{cn}^2(iu,\varkappa')$, und aus (10) folgt:
\[
\frac{dz}{\sqrt{z(1-z)(1-\varkappa^2z)}} = du,
\]
also
\[
z = \mathrm{sn}^2(u,\varkappa).
\]
Demnach ergibt (11)
\weq{12}{\mathrm{sn}\,(u,\varkappa) = \frac{-i\,\mathrm{sn}\,(iu,\varkappa')}{\mathrm{cn}\,(iu,\varkappa')}}
(worin das Vorzeichen aus dem speziellen Wert $u=0$ zu bestimmen ist).

Setzt man nun
\[
\int_0^{\frac{\pi}{2}} \frac{d\varphi}{\sqrt{1-\varkappa'^2\sin^2\varphi}} = K',
\]
so ergibt sich aus (8)
\[
\begin{aligned}
\mathrm{sn}\,(iu+2K',\varkappa') &= -\mathrm{sn}\,(iu,\varkappa'),\\
\mathrm{cn}\,(iu+2K',\varkappa') &= -\mathrm{cn}\,(iu,\varkappa'),
\end{aligned}
\]
und folglich aus (12)
\[
\mathrm{sn}\,(u-2iK',\varkappa) = \mathrm{sn}\,(u,\varkappa)
\]
%CONT>>
